\documentclass[12pt]{article}

\usepackage[margin=1in]{geometry}
\usepackage{amsmath,amssymb}
\usepackage{graphicx}
\usepackage{booktabs}
\usepackage{array}
\usepackage{siunitx}
\usepackage{enumitem}

\sisetup{per-mode=symbol}

\setlength{\parskip}{0.8ex}
\setlength{\parindent}{0pt}

\begin{document}

\title{Lab -- Module 3 -- PHYS 131\\[0.5em]Newton's Second Law}
\date{}
\maketitle

	extbf{Student Name:} {{ student_name }}\\ % FIELD:student_name:text
	extbf{Date:} {{ date }}\\ % FIELD:date:text
	extbf{Lab Section:} {{ section }}\\[1em] % FIELD:section:text

\section*{Introduction}

Newton's Second Law describes the relationship between net force, mass, and acceleration:
\[
\vec{F}_{\text{net}} = m \vec{a}
\]

In straight-line motion (1D motion), this reduces to:
\[
F_{\text{net}} = ma
\]

This experiment uses a cart on a track pulled by a hanging mass. The falling mass provides a nearly constant force, allowing you to measure the cart's acceleration and compare it with values predicted using Newton's Second Law.

\subsubsection*{Theoretical Model}

Consider a cart of mass $m_c$ on a horizontal track connected via a string to a hanging mass $m_h$. Neglecting friction:

\[
F_{\text{net}} = m_h g
\]

The total mass being accelerated is:
\[
m_{\text{total}} = m_c + m_h
\]

Thus, the expected (theoretical) acceleration is:
\[
a_{\text{theory}} = \frac{m_h g}{m_c + m_h}
\]

You will measure the actual (experimental) acceleration using timing data and compare the two.

\section*{Procedure}

\begin{enumerate}[label=\arabic*.]
    \item Measure the mass of the cart, $m_c$. Record this value in the Data Section.

    \item Measure the hanging mass $m_h$ for Trial Set 1 and record it.

    \item Set up the track so it is level. Attach the hanging mass to the cart using the string and pulley system.

    \item Pull the cart back slightly and release it from rest. Measure the time it takes to travel a known distance along the track. Repeat three times and record your three time measurements.

    \item Compute the cart's experimental acceleration for each trial using
    \[
    a_{\text{exp}} = \frac{2d}{t^2}
    \]

    \item Compute the theoretical acceleration using
    \[
    a_{\text{theory}} = \frac{m_h g}{m_c + m_h}
    \]

    \item Repeat Steps 2–6 for two additional hanging masses.

    \item Compare the experimental values to theoretical predictions and discuss possible sources of error.
\end{enumerate}

\section*{Data}

\subsection*{Cart Mass}
	extbf{Cart mass:} {{ cart.mass }} kg % FIELD:cart_mass:text

%========================
% TRIAL SET 1
%========================

\subsection*{Trial Set 1: Hanging Mass {{ set1.mh }} kg} % FIELD:set1_mh:text

\begin{center}
	extbf{Table 1}
\end{center}

% TABLE:set1_table:columns=d,t,aexp,atheory
\begin{center}
\begin{tabular}{cccc}
\toprule
\begin{tabular}{c}Distance\\$d$ (m)\end{tabular} &
\begin{tabular}{c}Time\\(s)\end{tabular} &
\begin{tabular}{c}Experimental\\Acceleration (m/s$^2$)\end{tabular} &
\begin{tabular}{c}Theoretical\\Acceleration (m/s$^2$)\end{tabular} \\
\midrule
{% for row in set1.rows %}
{{ row.d }} & {{ row.t }} & {{ row.aexp }} & {{ set1.atheory }} \\
{% endfor %}
\midrule
\textbf{Averages:} & {{ set1.avg.t }} & {{ set1.avg.aexp }} & {{ set1.atheory }} \\
\bottomrule
\end{tabular}
\end{center}

\vspace{1ex}

	extbf{Example calculation using average values:}\\ % FIELD:set1_example:textarea
{{ set1.example_calculation }}

%========================
% TRIAL SET 2
%========================

\subsection*{Trial Set 2: Hanging Mass {{ set2.mh }} kg} % FIELD:set2_mh:text

\begin{center}
	extbf{Table 2}
\end{center}

% TABLE:set2_table:columns=d,t,aexp,atheory
\begin{center}
\begin{tabular}{cccc}
\toprule
\begin{tabular}{c}Distance\\$d$ (m)\end{tabular} &
\begin{tabular}{c}Time\\(s)\end{tabular} &
\begin{tabular}{c}Experimental\\Acceleration (m/s$^2$)\end{tabular} &
\begin{tabular}{c}Theoretical\\Acceleration (m/s$^2$)\end{tabular} \\
\midrule
{% for row in set2.rows %}
{{ row.d }} & {{ row.t }} & {{ row.aexp }} & {{ set2.atheory }} \\
{% endfor %}
\midrule
\textbf{Averages:} & {{ set2.avg.t }} & {{ set2.avg.aexp }} & {{ set2.atheory }} \\
\bottomrule
\end{tabular}
\end{center}

%========================
% TRIAL SET 3
%========================

\subsection*{Trial Set 3: Hanging Mass {{ set3.mh }} kg} % FIELD:set3_mh:text

\begin{center}
	extbf{Table 3}
\end{center}

% TABLE:set3_table:columns=d,t,aexp,atheory
\begin{center}
\begin{tabular}{cccc}
\toprule
\begin{tabular}{c}Distance\\$d$ (m)\end{tabular} &
\begin{tabular}{c}Time\\(s)\end{tabular} &
\begin{tabular}{c}Experimental\\Acceleration (m/s$^2$)\end{tabular} &
\begin{tabular}{c}Theoretical\\Acceleration (m/s$^2$)\end{tabular} \\
\midrule
{% for row in set3.rows %}
{{ row.d }} & {{ row.t }} & {{ row.aexp }} & {{ set3.atheory }} \\
{% endfor %}
\midrule
\textbf{Averages:} & {{ set3.avg.t }} & {{ set3.avg.aexp }} & {{ set3.atheory }} \\
\bottomrule
\end{tabular}
\end{center}

\section*{Questions}

\begin{enumerate}[label=\arabic*.]
    \item How does changing the hanging mass affect the net force in the system?\\
    	extbf{Answer:} {{ q1 }} % FIELD:q1:textarea

    \item How does increasing the total mass of the system affect its acceleration?\\
    	extbf{Answer:} {{ q2 }} % FIELD:q2:textarea

    \item Compare the experimental and theoretical accelerations. Were they close? Explain possible differences.\\
    	extbf{Answer:} {{ q3 }} % FIELD:q3:textarea

    \item Identify two sources of experimental error in this setup.\\
    	extbf{Answer:} {{ q4 }} % FIELD:q4:textarea

    \item If friction were not negligible, how would this affect the acceleration? (Increase, decrease, or stay the same. Explain.)\\
    	extbf{Answer:} {{ q5 }} % FIELD:q5:textarea
\end{enumerate}

\end{document}
